\documentclass[a4paper,12pt]{article}

% ------------------------------------------------
% Кодировка, язык, математика
% ------------------------------------------------
\usepackage[T2A]{fontenc}
\usepackage[utf8]{inputenc}
\usepackage[russian]{babel}
\usepackage{amsmath,amssymb}
\usepackage{icomma}

% ------------------------------------------------
% Шрифт и типографика
% ------------------------------------------------
\usepackage{PTSans}
\renewcommand{\familydefault}{\sfdefault}
\usepackage{microtype}
\usepackage{setspace}
\onehalfspacing
\usepackage{parskip}

% ------------------------------------------------
% Страница
% ------------------------------------------------
\usepackage[
  a4paper,
  left=20mm,
  right=20mm,
  top=20mm,
  bottom=22mm,
  headheight=15pt
]{geometry}
\raggedbottom

% ------------------------------------------------
% Графика, таблицы, списки
% ------------------------------------------------
\usepackage{tikz}
\usepackage{tabularx,booktabs}
\usepackage{enumitem}

% ------------------------------------------------
% Цвета
% ------------------------------------------------
\usepackage{xcolor}
\definecolor{accent}{HTML}{3F7F7A}
\definecolor{darkaccent}{HTML}{285B57}
\definecolor{textgray}{HTML}{2F3437}
\definecolor{muted}{HTML}{687174}

% ------------------------------------------------
% Колонтитулы и PDF-метаданные
% ------------------------------------------------
\usepackage{fancyhdr}
\usepackage{hyperref}
\hypersetup{
  colorlinks=true,
  linkcolor=darkaccent,
  urlcolor=darkaccent,
  pdftitle={Чек-лист: теория вероятностей и статистика},
  pdfauthor={Лёвин Артём Александрович}
}

\pagestyle{fancy}
\fancyhf{}
\lhead{\small\textcolor{textgray}{Теория вероятностей и статистика}}
\rhead{\small\textcolor{textgray}{Ксения Клыкова}}
\cfoot{\small\thepage}
\renewcommand{\headrulewidth}{0.3pt}

% ------------------------------------------------
% Настройки
% ------------------------------------------------
\setlength{\parindent}{0pt}
\setlist[itemize]{leftmargin=6mm,itemsep=2pt,topsep=3pt}
\setlist[enumerate]{leftmargin=8mm,itemsep=4pt,topsep=4pt}

\newcommand{\key}[1]{\textcolor{darkaccent}{\textbf{#1}}}
\newcommand{\Me}{\operatorname{Me}}
\newcommand{\Mo}{\operatorname{Mo}}

\begin{document}
\pagenumbering{gobble}

% ================================================================
% ТИТУЛЬНЫЙ ЛИСТ
% ================================================================
\begin{titlepage}
\thispagestyle{empty}

\begin{center}
\vspace*{24mm}

{\Large\textcolor{darkaccent}{\textbf{ЧЕК-ЛИСТ}}}

\vspace{9mm}

{\huge\textbf{Теория вероятностей и статистика}}

\vspace{8mm}

{\Large
Выбор без возвращения, характеристики центра,\\[2mm]
размах, дисперсия и стандартное отклонение
}

\vspace{12mm}

{\large Подготовка к ЕГЭ}

\vfill

{\large
\textbf{Лёвин Артём Александрович}\\[2mm]
эксклюзивно для Ксении Клыковой
}

\vspace{8mm}

{\large \today}

\vfill

\begin{tikzpicture}[remember picture,overlay]
\draw[color=accent,line width=1.1pt]
([xshift=18mm,yshift=17mm]current page.south west)
.. controls
([xshift=58mm,yshift=29mm]current page.south west)
and
([xshift=-63mm,yshift=8mm]current page.south east)
..
([xshift=-18mm,yshift=18mm]current page.south east);
\end{tikzpicture}

\end{center}
\end{titlepage}
\pagenumbering{arabic}

% ================================================================
\section*{1. Вероятность: сначала переводим текст в последовательность событий}

В задачах с последовательными извлечениями главная трудность часто состоит в правильной интерпретации условия. Перед вычислениями полезно записать, \key{что должно произойти на каждом шаге}.

Например, фраза

\begin{center}
\textit{«первый синий фломастер появился третьим»}
\end{center}

означает, что первые два фломастера должны быть несиними, а третий --- синим. Если фломастеры только красные и синие, схема события такова:

\[
K\;K\;S,
\]

где $K$ --- красный, $S$ --- синий.

\subsection*{Выбор без возвращения}

Если извлечённый предмет обратно не возвращают, после каждого шага меняются и общее количество предметов, и количество предметов соответствующего цвета.

Пусть имеется $7$ красных и $3$ синих фломастера. Тогда

\[
P(KKS)
=
\frac{7}{10}\cdot\frac{6}{9}\cdot\frac{3}{8}
=
\frac{7}{40}.
\]

Знаменатели уменьшаются:

\[
10,\quad 9,\quad 8,
\]

поскольку после каждого извлечения предметов становится на один меньше.

\subsection*{Правило произведения}

Для конкретной последовательности событий используем условные вероятности:

\[
P(A\cap B\cap C)
=
P(A)\,P(B\mid A)\,P(C\mid A\cap B).
\]

Практический смысл: если требуется, чтобы произошло событие $A$ \textbf{и затем} $B$ \textbf{и затем} $C$, вероятности по ветви перемножаются.

% ================================================================
\section*{2. Когда вероятности складывают}

Если условию подходят несколько \key{взаимоисключающих сценариев}, их вероятности складывают.

Например, две пары «чашка + блюдце» могут оказаться одноцветными в трёх случаях по количеству синих пар:

\[
0\text{ синих пар},\qquad 1\text{ синяя пара},\qquad 2\text{ синие пары}.
\]

Эти случаи одновременно произойти не могут, поэтому

\[
P=P_0+P_1+P_2.
\]

Внутри каждого сценария последовательные извлечения по-прежнему требуют умножения вероятностей.

\begin{center}
\renewcommand{\arraystretch}{1.3}
\begin{tabularx}{0.9\linewidth}{@{}lX@{}}
\toprule
\textbf{Фраза в условии} & \textbf{Что делаем} \\
\midrule
$A$ и затем $B$ & перемножаем вероятности по одной ветви \\
$A$ или $B$, причём случаи несовместимы & складываем вероятности случаев \\
\bottomrule
\end{tabularx}
\end{center}

% ================================================================
\section*{3. Задача о чашках и блюдцах}

На первой полке стоят $36$ блюдец: $14$ синих и $22$ красных. На второй полке стоят $36$ чашек: $27$ синих и $9$ красных. С каждой полки случайно выбирают по два предмета без возвращения.

Чтобы из выбранных предметов можно было составить две одноцветные чайные пары, число выбранных синих чашек должно совпасть с числом выбранных синих блюдец. Возможны три случая: выбрать по $0$, по $1$ или по $2$ синих предмета каждого вида.

Для блюдец:

\[
P_{\text{бл}}(0)
=
\frac{22}{36}\cdot\frac{21}{35}
=
\frac{11}{30},
\]

\[
P_{\text{бл}}(1)
=
\frac{14}{36}\cdot\frac{22}{35}
+
\frac{22}{36}\cdot\frac{14}{35}
=
\frac{22}{45},
\]

\[
P_{\text{бл}}(2)
=
\frac{14}{36}\cdot\frac{13}{35}
=
\frac{13}{90}.
\]

Для чашек:

\[
P_{\text{ч}}(0)=\frac{9}{36}\cdot\frac{8}{35}=\frac{2}{35},
\qquad
P_{\text{ч}}(1)=\frac{27}{70},
\qquad
P_{\text{ч}}(2)=\frac{39}{70}.
\]

Следовательно,

\[
\begin{aligned}
P
&=
P_{\text{бл}}(0)P_{\text{ч}}(0)
+P_{\text{бл}}(1)P_{\text{ч}}(1)
+P_{\text{бл}}(2)P_{\text{ч}}(2)\\
&=
\frac{11}{30}\cdot\frac{2}{35}
+
\frac{22}{45}\cdot\frac{27}{70}
+
\frac{13}{90}\cdot\frac{39}{70}
=
\frac{29}{100}.
\end{aligned}
\]

\key{Самопроверка:} для каждого вида предметов вероятности случаев $0$, $1$, $2$ синих должны в сумме давать $1$.

% ================================================================
\section*{4. Характеристики центра данных}

На занятии использовались три основные характеристики:

\[
\boxed{\text{среднее арифметическое}}
\qquad
\boxed{\text{медиана}}
\qquad
\boxed{\text{мода}}.
\]

Они отвечают на разные вопросы, поэтому выбор характеристики зависит от смысла данных.

\subsection*{4.1. Среднее арифметическое}

Для числового набора $x_1,x_2,\ldots,x_n$

\[
\overline{x}
=
\frac{x_1+x_2+\cdots+x_n}{n}.
\]

Среднее использует каждое наблюдение и поэтому чувствительно к очень большим или очень маленьким значениям.

Например, для набора

\[
0{,}1,\quad 1,\quad 2,\quad 3,\quad 100
\]

имеем

\[
\overline{x}
=
\frac{106{,}1}{5}
=
21{,}22.
\]

Одно большое значение $100$ заметно сдвинуло среднее вверх.

\subsection*{4.2. Медиана}

\key{Медиана} --- значение, характеризующее середину \textbf{упорядоченного} числового набора.

Первый шаг всегда один:

\[
\boxed{\text{расположить данные по возрастанию}}.
\]

Если количество элементов нечётное, медиана --- центральный элемент. Например,

\[
3,\ 4,\ 5,\ 8,\ 10
\qquad\Longrightarrow\qquad
\Me=5.
\]

Если количество элементов чётное, медиана равна среднему арифметическому двух центральных элементов:

\[
2,\ 4,\ 5,\ 6,\ 7,\ 9
\qquad\Longrightarrow\qquad
\Me=\frac{5+6}{2}=5{,}5.
\]

Для набора

\[
0{,}1,\quad 1,\quad 2,\quad 3,\quad 100
\]

медиана равна $2$. Она значительно устойчивее среднего к отдельному большому значению $100$.

\subsection*{4.3. Мода}

\key{Мода} --- значение, встречающееся чаще остальных.

Например,

\[
2,\ 3,\ 3,\ 4,\ 4,\ 4,\ 5
\qquad\Longrightarrow\qquad
\Mo=4.
\]

Мода --- само значение $4$, а число $3$ показывает его абсолютную частоту.

Если максимальная частота одинакова у нескольких значений, мод может быть несколько. Например,

\[
1,\ 1,\ 2,\ 2,\ 3
\]

имеет две моды: $1$ и $2$.

% ================================================================
\section*{5. Как выбрать характеристику}

\begin{center}
\renewcommand{\arraystretch}{1.35}
\begin{tabularx}{\linewidth}{@{}lXX@{}}
\toprule
\textbf{Характеристика} & \textbf{Что учитывает} & \textbf{Когда особенно полезна} \\
\midrule
Среднее арифметическое & все числовые наблюдения & когда важен вклад каждого значения; чувствительность к выбросам нужно учитывать \\
Медиана & порядок данных и центральное положение & при сильных выбросах или заметной асимметрии \\
Мода & максимальную частоту & когда важно определить самое частое значение или категорию \\
\bottomrule
\end{tabularx}
\end{center}

\subsection*{Пример: время дороги до школы}

Пусть почти все ученики добираются до школы за $10$--$20$ минут, а один ученик --- за $60$ минут. Значение $60$ является выбросом и заметно увеличивает среднее. Для описания типичного времени в такой ситуации медиана обычно устойчивее.

\subsection*{Числовые и категориальные данные}

Пусть способы дороги до школы имеют частоты:

\begin{center}
\begin{tabular}{@{}lrrrr@{}}
\toprule
\textbf{Способ} & машина & метро & пешком & автобус \\
\midrule
\textbf{Частота} & 4 & 7 & 8 & 11 \\
\bottomrule
\end{tabular}
\end{center}

Среднее чисел $4,7,8,11$ равно $7{,}5$ и означает среднее количество учеников на одну категорию. Оно не отвечает на вопрос, \textit{какой способ передвижения встречается чаще всего}.

Для категориальных данных без естественного порядка медиана не определяется. Здесь содержательно полезна мода: наиболее частая категория --- \key{автобус}.

% ================================================================
\section*{6. Разброс данных}

Характеристики центра показывают, около чего расположены данные. Чтобы описать, насколько далеко значения разбросаны друг от друга, используют характеристики разброса.

\subsection*{6.1. Размах}

Размах равен разности максимального и минимального значений:

\[
R=x_{\max}-x_{\min}.
\]

Например, для набора

\[
1,\ 2,\ 3,\ 3,\ 4,\ 5,\ 5,\ 5,\ 10
\]

получаем

\[
R=10-1=9.
\]

Размах учитывает только два крайних значения, поэтому даёт грубую оценку разброса.

\subsection*{6.2. Дисперсия}

Пусть $x_1,x_2,\ldots,x_n$ --- весь рассматриваемый числовой набор, а $\overline{x}$ --- его среднее арифметическое. Дисперсия этого набора определяется формулой

\[
D
=
\frac{1}{n}
\sum_{i=1}^{n}(x_i-\overline{x})^2.
\]

В развёрнутом виде:

\[
D
=
\frac{
(x_1-\overline{x})^2+
(x_2-\overline{x})^2+
\cdots+
(x_n-\overline{x})^2
}{n}.
\]

\subsection*{Алгоритм вычисления}

\begin{enumerate}
    \item Найдите среднее арифметическое $\overline{x}$.
    \item Для каждого элемента найдите отклонение $x_i-\overline{x}$.
    \item Возведите каждое отклонение в квадрат.
    \item Сложите квадраты отклонений.
    \item Разделите сумму на количество наблюдений $n$.
\end{enumerate}

Квадраты не дают положительным и отрицательным отклонениям взаимно уничтожиться при сложении.

Если значение $a$ встречается $k$ раз, одинаковые слагаемые удобно объединять:

\[
\underbrace{(a-\overline{x})^2+\cdots+(a-\overline{x})^2}_{k\ \text{раз}}
=
k(a-\overline{x})^2.
\]

\subsection*{6.3. Стандартное отклонение}

Стандартное отклонение равно квадратному корню из дисперсии:

\[
\sigma=\sqrt{D}.
\]

Дисперсия измеряется в квадрате исходных единиц, а стандартное отклонение --- в тех же единицах, что и сами данные. Поэтому $\sigma$ удобнее для интерпретации масштаба разброса.

Чем больше $\sigma$, тем сильнее значения в среднем удалены от среднего арифметического. Если все элементы набора одинаковы, то

\[
D=0,
\qquad
\sigma=0.
\]

Интервал

\[
[\overline{x}-\sigma;\ \overline{x}+\sigma]
\]

можно использовать как наглядный ориентир масштаба одного стандартного отклонения от среднего. Без дополнительных сведений о распределении данных этому интервалу нельзя приписывать фиксированную вероятность попадания.

% ================================================================
\section*{7. Полный пример: дисперсия и стандартное отклонение}

Рассмотрим набор

\[
1,\ 2,\ 3,\ 3,\ 4,\ 5,\ 5,\ 5,\ 10,\ 12.
\]

Количество элементов равно $10$, а сумма равна $50$, поэтому

\[
\overline{x}=\frac{50}{10}=5.
\]

Дисперсия:

\[
\begin{aligned}
D
&=
\frac{
(1-5)^2+(2-5)^2+2(3-5)^2+(4-5)^2
+3(5-5)^2+(10-5)^2+(12-5)^2
}{10}
\\[1mm]
&=
\frac{16+9+8+1+0+25+49}{10}
=
\frac{108}{10}
=
10{,}8.
\end{aligned}
\]

Стандартное отклонение:

\[
\sigma=\sqrt{10{,}8}\approx3{,}29.
\]

Границы интервала одного стандартного отклонения от среднего:

\[
5-3{,}29\approx1{,}71,
\qquad
5+3{,}29\approx8{,}29.
\]

То есть

\[
[\overline{x}-\sigma;\ \overline{x}+\sigma]
\approx
[1{,}71;\ 8{,}29].
\]

Этот интервал описывает масштаб разброса около среднего; сам по себе он не гарантирует заданную долю наблюдений внутри него.

% ================================================================
\section*{8. Что происходит при выбросе}

Рассмотрим набор

\[
2,\ 4,\ 5,\ 7,\ 9,\ 11,\ 13.
\]

Его медиана равна

\[
\Me=7.
\]

Заменим $13$ на $30$:

\[
2,\ 4,\ 5,\ 7,\ 9,\ 11,\ 30.
\]

Медиана по-прежнему равна $7$, а среднее арифметическое увеличится с

\[
\frac{51}{7}\approx7{,}29
\]

до

\[
\frac{68}{7}\approx9{,}71.
\]

Это показывает устойчивость медианы к изменению крайнего значения.

\key{Важно:} медиана зависит от центрального положения элементов после упорядочивания. Она может измениться и при неизменном количестве элементов, если изменятся значения в центральной части ряда.

% ================================================================
\section*{9. Частые ошибки}

\begin{enumerate}
    \item \key{Искать медиану до упорядочивания данных.} Сначала расположите числа по возрастанию.

    \item \key{При чётном количестве элементов выбирать только одно центральное число.} Нужно взять среднее двух центральных значений.

    \item \key{Путать моду и частоту.} Мода --- значение или категория; частота --- количество появлений.

    \item \key{При выборе без возвращения оставлять прежний знаменатель.} После каждого извлечения общее количество оставшихся предметов уменьшается.

    \item \key{Складывать вероятности событий, которые должны произойти последовательно.} По одной последовательной ветви вероятности перемножаются.

    \item \key{Забывать оба порядка в смешанном случае.} Например, «один синий и один красный» может произойти как $SK$ или как $KS$.

    \item \key{Механически применять среднее, медиану и моду к любым данным.} Сначала определите, являются данные числовыми или категориальными и какой вопрос требуется решить.

    \item \key{Считать интервал $\overline{x}\pm\sigma$ гарантированным вероятностным интервалом.} Без дополнительных предпосылок это только характеристика масштаба разброса.
\end{enumerate}

% ================================================================
\section*{10. Итоговый чек-лист}

\begin{itemize}
    \item[$\square$] Я умею переводить словесное условие вероятностной задачи в последовательность событий.
    \item[$\square$] Я учитываю изменение числителей и знаменателей при выборе без возвращения.
    \item[$\square$] Я различаю ситуации «и» и «или».
    \item[$\square$] Я умею находить среднее арифметическое, медиану и моду.
    \item[$\square$] Перед поиском медианы я упорядочиваю числовой набор.
    \item[$\square$] Я понимаю влияние выброса на среднее и медиану.
    \item[$\square$] Я умею находить размах.
    \item[$\square$] Я умею вычислять дисперсию и стандартное отклонение.
    \item[$\square$] Я могу объяснить смысл полученной статистической характеристики.
\end{itemize}

% ================================================================
% ТРЕНИРОВОЧНЫЙ БЛОК
% ================================================================
\newpage
\section*{Тренировочный блок}

\textcolor{muted}{Сначала определите, какое правило требуется применить. В вероятностных задачах выпишите последовательность событий; в статистических --- укажите смысл найденной характеристики.}

\begin{enumerate}[label=\textbf{\arabic*.},leftmargin=9mm,itemsep=4mm]

\item \textbf{Характеристики центра и размах.}

Дан набор
\[
12,\quad 5,\quad 8,\quad 5,\quad 17,\quad 9,\quad 5,\quad 6.
\]
Найдите среднее арифметическое, медиану, моду и размах.

\item \textbf{Первый синий --- третьим.}

В коробке находятся $8$ красных и $4$ синих фломастера. Фломастеры последовательно извлекают случайным образом без возвращения. Найдите вероятность того, что первый синий фломастер появится третьим.

\item \textbf{Выброс и типичное значение.}

Время дороги до школы у семи учеников составило
\[
12,\quad14,\quad15,\quad15,\quad16,\quad18,\quad65
\]
минут.

\begin{enumerate}[label=\alph*),leftmargin=7mm,itemsep=2pt,topsep=2pt]
    \item Найдите среднее арифметическое.
    \item Найдите медиану.
    \item Какая из этих двух характеристик устойчивее описывает типичное время дороги для данного набора? Кратко объясните выбор.
\end{enumerate}

\newpage

\item \textbf{Дисперсия и стандартное отклонение.}

Для набора
\[
1,\quad3,\quad3,\quad5,\quad8
\]
найдите среднее арифметическое, дисперсию
\[
D=\frac1n\sum_{i=1}^{n}(x_i-\overline{x})^2,
\]
стандартное отклонение $\sigma$ и границы интервала $[\overline{x}-\sigma;\ \overline{x}+\sigma]$. Ответы с корнем округлите до сотых.

\item \textbf{Чашки и блюдца.}

На одной полке находится $10$ блюдец: $6$ синих и $4$ красных. На другой полке находится $10$ чашек: $7$ синих и $3$ красных.

С каждой полки случайным образом без возвращения выбирают по два предмета. Найдите вероятность того, что выбранные предметы можно разбить на две пары «чашка + блюдце», причём в каждой паре чашка и блюдце одного цвета.

\end{enumerate}

% ================================================================
% ОТВЕТЫ
% ================================================================
\newpage
\section*{Ответы}

\renewcommand{\arraystretch}{1.45}
\begin{center}
\begin{tabularx}{\linewidth}{@{}p{12mm}X@{}}
\toprule
\textbf{№} & \textbf{Ответ} \\
\midrule
1 & $\overline{x}=8{,}375$; $\Me=7$; $\Mo=5$; $R=12$. \\
2 & $\displaystyle \frac{28}{165}\approx0{,}170$. \\
3 & $\displaystyle \overline{x}=\frac{155}{7}\approx22{,}14$ мин; $\Me=15$ мин. Для описания типичного времени в этом наборе устойчивее медиана: значение $65$ мин является сильным выбросом. \\
4 & $\overline{x}=4$; $D=\dfrac{28}{5}=5{,}6$; $\sigma=\sqrt{5{,}6}\approx2{,}37$; интервал $\approx[1{,}63;\ 6{,}37]$. \\
5 & $\displaystyle \frac{31}{75}\approx0{,}413$. \\
\bottomrule
\end{tabularx}
\end{center}

\end{document}